\section{Собственные числа и собственные векторы линейного оператора. Теорема о нахождении собственных чисел. Спектр оператора. Нахождение собственных векторов. Основное свойство собственных векторов}

\begin{definition}
Ненулевой вектор $x \in V$ ($x \neq 0$) называется \textbf{собственным вектором} линейного оператора $\varphi \in L(V, V)$, если при действии оператора он переходит в вектор, коллинеарный самому себе, то есть:
\[
\varphi(x) = \lambda x,
\]
где $\lambda \in F$ --- некоторое число, называемое \textbf{собственным числом} (или собственным значением), соответствующим этому собственному вектору.
\end{definition}

\begin{theorem}[О нахождении собственных чисел]
Число $\lambda$ является собственным числом оператора $\varphi$ тогда и только тогда, когда оно является корнем характеристического уравнения:
\[
\det(A_\varphi - \lambda E_n) = 0.
\]
\end{theorem}

\begin{proof}
Пусть $x \neq 0$ --- собственный вектор, соответствующий числу $\lambda$, и $X$ --- его столбец координат в базисе $E$. Равенство $\varphi(x) = \lambda x$ в координатной форме записывается как:
\[
A_\varphi X = \lambda X \quad \Rightarrow \quad (A_\varphi - \lambda E_n) X = 0.
\]
Это однородная система линейных алгебраических уравнений относительно неизвестных координат вектора $X$. Так как $x \neq 0$, система должна иметь нетривиальное решение ($X \neq 0$). Известно, что однородная система имеет нетривиальное решение тогда и только тогда, когда определитель её матрицы коэффициентов равен нулю. Следовательно, $\det(A_\varphi - \lambda E_n) = 0$.
Обратное утверждение доказывается аналогично: если определитель равен нулю, система имеет нетривиальное решение $X \neq 0$, которое задает координаты собственного вектора $x$.
\hfill$\blacksquare$
\end{proof}

\begin{definition}
Множество всех собственных чисел линейного оператора $\varphi$ называется его \textbf{спектром}.
\end{definition}

\begin{remark}[Нахождение собственных векторов]
Для нахождения собственных векторов, соответствующих данному собственному числу $\lambda$, необходимо:
\begin{enumerate}
    \item Составить матрицу $A_\varphi - \lambda E_n$.
    \item Найти фундаментальную систему решений (ФСР) однородной системы $(A_\varphi - \lambda E_n)X = 0$.
    \item Любая ненулевая линейная комбинация векторов этой ФСР будет задавать координаты собственного вектора, соответствующего $\lambda$.
\end{enumerate}
\end{remark}

\begin{theorem}[Основное свойство собственных векторов]
Собственные векторы $x_1, x_2, \dots, x_k$, соответствующие попарно различным собственным числам $\lambda_1, \lambda_2, \dots, \lambda_k$ линейного оператора $\varphi$, линейно независимы.
\end{theorem}

\begin{proof}
Доказательство проведем методом математической индукции по числу векторов $k$.
\begin{enumerate}
    \item \textit{База индукции ($k=1$):} Вектор $x_1$ является собственным, следовательно, по определению $x_1 \neq 0$. Система из одного ненулевого вектора линейно независима.
    \item \textit{Индукционный переход:} Предположим, что утверждение верно для любых $k$ собственных векторов с различными собственными числами. Рассмотрим $k+1$ вектор $x_1, \dots, x_{k+1}$ с собственными числами $\lambda_1, \dots, \lambda_{k+1}$ (все $\lambda_i$ попарно различны).
    Составим линейную комбинацию, равную нулю:
    \[
    c_1 x_1 + c_2 x_2 + \dots + c_k x_k + c_{k+1} x_{k+1} = 0. \quad (*)
    \]
    Применим к обеим частям оператор $\varphi$. Используя свойство $\varphi(x_i) = \lambda_i x_i$, получим:
    \[
    c_1 \lambda_1 x_1 + c_2 \lambda_2 x_2 + \dots + c_k \lambda_k x_k + c_{k+1} \lambda_{k+1} x_{k+1} = 0. \quad (**)
    \]
    Умножим равенство $(*)$ на $\lambda_{k+1}$ и вычтем его из равенства $(**)$:
    \[
    c_1 (\lambda_1 - \lambda_{k+1}) x_1 + c_2 (\lambda_2 - \lambda_{k+1}) x_2 + \dots + c_k (\lambda_k - \lambda_{k+1}) x_k = 0.
    \]
    Полученная линейная комбинация содержит только $k$ векторов $x_1, \dots, x_k$. По предположению индукции они линейно независимы, поэтому все коэффициенты при них равны нулю:
    \[
    c_i (\lambda_i - \lambda_{k+1}) = 0 \quad \text{для } i = 1, \dots, k.
    \]
    Так как собственные числа попарно различны, $\lambda_i - \lambda_{k+1} \neq 0$, откуда следует $c_i = 0$ для всех $i = 1, \dots, k$.
    Подставляя $c_1 = \dots = c_k = 0$ в исходное равенство $(*)$, получаем $c_{k+1} x_{k+1} = 0$. Поскольку $x_{k+1} \neq 0$, имеем $c_{k+1} = 0$.
\end{enumerate}
Таким образом, все коэффициенты $c_i$ равны нулю, что доказывает линейную независимость системы $x_1, \dots, x_{k+1}$.
\hfill$\blacksquare$
\end{proof}
